% Guard: macros.tex is \input both by main.tex (after the packages it needs,
% e.g. amsmath for \DeclareMathOperator) and by individual section files so
% they remain compilable on their own. Loading it twice would redefine every
% \newcommand and abort, so the second and later loads are no-ops.
\ifdefined\PQRmacrosloaded\endinput\fi
\let\PQRmacrosloaded\relax

% --- 0. revision marking ---
% Text added to Section 2 during the current revision is wrapped in \added{...}
% so the advisors can see at a glance what is new. Redefine to \relax to print
% the paper without the marking.
\newcommand{\added}[1]{\textcolor{blue}{#1}}
\long\def\added#1{\begingroup\color{blue}#1\endgroup}  % \long: may span paragraphs


% --- A. operator-style score and function names ---
\newcommand{\HD}{\mathit{HD}}                        % histogram distance; use \HD_\gamma
\newcommand{\Attr}{\mathit{Attr}}                    % attribute set of a schema
\newcommand{\Attrs}{\mathit{Attrs}}                  % attribute set argument
% Abbreviated in MATH only; the prose still spells the names out in full.
% Spelling them out inside formulas makes displayed equations far too wide
% for a two-column layout (the product-rule step below overflowed by 200pt).
\newcommand{\Use}{\mathit{Use}}
\newcommand{\Cov}{\mathit{Cov}}
\newcommand{\Rep}{\mathit{Rep}}
\newcommand{\DivS}{\mathit{Div}}
\newcommand{\TVD}{\mathit{TVD}}
\newcommand{\GS}{\mathit{GS}}                        % global sensitivity
\DeclareMathOperator{\Dom}{Dom}
\DeclareMathOperator{\Lap}{Laplace}

% --- B. core objects ---
\newcommand{\Qone}{Q_1}
\newcommand{\Qtwo}{Q_2}
\newcommand{\pD}{p(D)}
\newcommand{\views}{\mathcal{P}}
\newcommand{\hist}[2]{H_{#1,#2}}                     % \hist{\Qone(\pD)}{A}
\newcommand{\dhist}{\Delta H}
\newcommand{\nhist}{\widetilde{\Delta H}}            % the released noisy object

% --- C. privacy budget and parameters ---
\newcommand{\eps}{\epsilon}
\newcommand{\epsP}{\epsilon_P}
\newcommand{\epsS}{\epsilon_S}
\newcommand{\epsH}{\epsilon_H}
\newcommand{\gammaref}{\gamma_{\mathrm{ref}}}
\newcommand{\gammaeval}{\gamma_{\mathrm{eval}}}
\newcommand{\taumax}{\tau_{\max}}

% --- D. mechanisms ---
\newcommand{\ExpMech}{\textsc{ExpMech}}
\newcommand{\OneShotTop}[1]{\textsc{OneShotTop-}#1}
\newcommand{\AdaptiveTau}{\textsc{AdaptiveTau}}
% --- E. experiment baselines (used in sections/experiments.tex) ---
\newcommand{\Bone}{\ensuremath{\mathsf{B1}}}
\newcommand{\Btwo}{\ensuremath{\mathsf{B2}}}
\newcommand{\Bthree}{\ensuremath{\mathsf{B3}}}

\section{Preliminaries}
\label{sec:preliminaries}

In this work, we study the problem of generating differentially private explanations
for semantic differences between query outputs.
This section introduces the formal concepts and notation used throughout the paper.

\subsection{Databases, Queries, and Predicate Views}
\label{subsec:prelim-db-queries}

\paragraph{Databases.}
Let $D$ be a relational database over schema $\mathcal{R}$~\cite{Foundations}.
Each tuple $t \in D$ corresponds to an individual entity or record, and we write
$N=|D|$ for the number of tuples.
We denote by $\Attr(D)=\{A_1,\ldots,A_m\}$ the set of attributes appearing in the
schema of $D$.
Each attribute $A$ has a domain $\Dom(A)$, which may be categorical or numerical,
and each tuple $t\in D$ assigns a value $t[A]\in\Dom(A)$ to every attribute $A$.
Throughout, $\Dom(A)$ is the domain declared by the schema and is therefore
\emph{public}, since it does not depend on the contents of $D$.
This distinction matters once privacy enters the picture, since a domain read off
the data at hand would itself disclose information about the data
(Section~\ref{subsec:prelim-dp}).

\paragraph{Queries.}
We consider SQL-style selection and projection queries together with aggregate
queries, all of them having Boolean selection predicates over a single
relational table $D$.
A query $Q$ consists of a selection predicate $\varphi$ over $\Attr(D)$, and
aggregate queries have the form
\[
\texttt{SELECT } agg(A)
\texttt{ FROM } D
\texttt{ WHERE } \varphi ,
\]
where $agg \in \{\texttt{COUNT},\texttt{SUM},\texttt{AVG}\}$ and
$\varphi$ is a Boolean combination of \emph{atomic conditions}, namely numerical
comparisons such as $A \leq c$ or $A \geq c$, categorical equality predicates
such as $A = v$, and categorical set-membership predicates such as $A \in S$.

Our explanations describe the \emph{population} a query selects rather than the
scalar it may return, since it is the population whose composition changes when a
query is refined or a database is repaired, and an aggregate value alone admits no
distributional summary.
We therefore write
\[
Q(D) \;=\; \sigma_{\varphi}(D) \;=\; \{t \in D \mid \varphi(t)=\texttt{true}\}
\]
for the \emph{output} of $Q$ on $D$, that is, the set of tuples selected by $Q$.
For an aggregate query the reported value is a function of $Q(D)$, and all
statistics defined below are taken over $Q(D)$.

\paragraph{Predicate views.}
To explain semantic differences between query outputs, our framework analyzes
subsets of the data defined by predicates.
Predicate-defined subpopulations are commonly used in explanation and
query-analysis frameworks to localize meaningful changes in the
data~\cite{FEDEX,TaoGMR22,wu2013}.

\begin{definition}[Predicate View]
\label{def:predicate-view}
Given a database $D$ and a predicate $p$, the predicate view induced by $p$ is
\[
p(D)=\{t\in D \mid p(t)=\texttt{true}\}.
\]
\end{definition}

The framework is given a finite collection $\views$ of view predicates, which
localize the explanations to specific semantic regions of the data.
We take $\views$ to consist of conjunctions of at most $\ell$ atomic conditions over
a designated set of view attributes, with $\ell$ small (in our experiments
$\ell\le 2$); its members are enumerated from the public schema and domain
metadata, so both $\views$ and its size $|\views|$ are fixed independently of the
contents of $D$.
Two further public attribute sets are used.
The \emph{explanation attributes} $\Attrs\subseteq\Attr(D)$ are those whose
distributions may be reported, and the \emph{reference attributes}
$\mathcal{A}_R\subseteq\Attr(D)$ are those against which a view is compared to the
full database in Section~\ref{sec:explanation-quality}.

Evaluating a query inside a view is the basic operation of the framework, and the
following identity is what makes it well behaved.

\begin{proposition}[View restriction]
\label{prop:view-restriction}
For every query $Q$ with selection predicate $\varphi$ and every view predicate $p$,
\[
Q(p(D)) \;=\; \sigma_{\varphi\wedge p}(D) \;=\; Q(D)\cap p(D).
\]
In particular $Q(p(D)) \subseteq p(D)$ and $Q(p(D)) \subseteq Q(D)$.
\end{proposition}

\begin{proof}
Both $Q$ and $p$ act on the same single relation by selection, and selections
compose, so that $t\in\sigma_\varphi(\sigma_p(D))$ holds iff $t\in D$, $p(t)$ and $\varphi(t)$,
which is the definition of $\sigma_{\varphi\wedge p}(D)$, and also of
$\sigma_\varphi(D)\cap\sigma_p(D)$.
The two containments are immediate.
\end{proof}

The containment $Q(p(D))\subseteq p(D)$ is used repeatedly.
It is why the histogram
of a query output inside a view is supported on the bins of the view itself
(Definition~\ref{def:hd}), and why coverage inside a view is bounded by the view
size (Section~\ref{sec:explanation-quality}).

Throughout the paper we study pairs of related queries $(Q_1,Q_2)$ evaluated over
the same database $D$, and our goal is to explain the semantic differences between
the outputs $Q_1(D)$ and $Q_2(D)$.
The next subsection describes two application settings that naturally give rise to
such pairs.

\begin{example}[Running example]
\label{ex:running}
We use the Adult refinement of Section~\ref{sec:introduction} as a running example.
The database $D$ holds census records over attributes including
\texttt{gender}, \texttt{education}, \texttt{occupation}, \texttt{hoursPerWeek} and
\texttt{income}.
The original query $\Qone$ and its refinement $\Qtwo$ have selection
predicates
\begin{align*}
\varphi_1 &\;=\; \texttt{education} \in \{\texttt{Bachelors},\texttt{Masters}\}\\
          &\hphantom{\;=\;} \wedge\; \texttt{hoursPerWeek} \ge 50 ,\\[3pt]
\varphi_2 &\;=\; \texttt{education} = \texttt{Masters}
              \;\wedge\; \texttt{hoursPerWeek} \ge 30 .
\end{align*}
so that $\Qone(D)$ and $\Qtwo(D)$ are the two populations to be compared.
A view predicate such as $p = (\texttt{education}=\texttt{Masters})$ isolates a
semantic region of $D$, and by Proposition~\ref{prop:view-restriction},
$\Qone(\pD)$ is exactly the set of Master's-degree holders working at least
$50$ hours per week.
\end{example}

\subsection{Query Refinement and Consistent Query Answering}
\label{subsec:prelim-settings}

Both settings we consider derive a second output from an original query, and it is
the resulting pair that our framework takes as input. We state them formally.

\paragraph{Query refinement.}
Refinement modifies a query so that its output meets a set of constraints on the
composition of the result while staying close to the original query
\cite{MK09,LMSJ23,ERICA,CSSM24}.

\begin{definition}[Output constraint]
\label{def:output-constraint}
Let $G\subseteq\Attr(D)$ be a set of grouping attributes and let $g$ be an
assignment of values to $G$, identifying a group of tuples. A \emph{cardinality
constraint} is a triple $\langle g,\alpha,\beta\rangle$, satisfied by a set of
tuples $T$ when
\[
\alpha \;\le\; |\{t\in T \mid t[G]=g\}| \;\le\; \beta ,
\]
and a \emph{proportional constraint} replaces this count by its fraction of $|T|$.
For a set $\mathcal{C}$ of such constraints we write $T\models\mathcal{C}$ when $T$
satisfies every member of $\mathcal{C}$.
\end{definition}

\begin{definition}[Query refinement]
\label{def:refinement}
Let $Q$ have selection predicate $\varphi$, and let $\Phi_\varphi$ be the space of
predicates obtained from $\varphi$ by relaxing, tightening, removing or replacing
its atomic conditions. Given $D$, a constraint set $\mathcal{C}$ and a deviation
measure $d$ on predicates, a \emph{refinement} of $Q$ is a query $Q'$ with
predicate $\varphi'\in\Phi_\varphi$ such that $Q'(D)\models\mathcal{C}$. A
refinement mechanism returns such a $Q'$, typically one minimizing
$d(\varphi,\varphi')$.
\end{definition}

\paragraph{Consistent query answering.}
CQA evaluates a query over a database that may violate a set of integrity
constraints, by reference to the instances that repair those violations
\cite{arenas1999,bertossi2006,bertossi2011}.

\begin{definition}[Repair]
\label{def:repair}
Let $\Sigma$ be a set of integrity constraints. An instance $D'$ is a
\emph{repair} of $D$ with respect to $\Sigma$ if $D'\subseteq D$,
$D'\models\Sigma$, and no proper superset of $D'$ contained in $D$ satisfies
$\Sigma$. We write $\mathsf{Repairs}(D,\Sigma)$ for the set of all repairs of $D$.
\end{definition}

\begin{definition}[Consistent answer]
\label{def:consistent-answer}
The \emph{consistent answer} to a query $Q$ over $D$ with respect to $\Sigma$ is
\[
Q_\Sigma(D)=\bigcap_{D'\in \mathsf{Repairs}(D,\Sigma)} Q(D') ,
\]
the set of tuples returned by $Q$ on every repair of $D$.
\end{definition}

We take repairs to be subsets of $D$, as in Definition~\ref{def:repair}, so that
$Q_\Sigma(D)\subseteq Q(D)\subseteq D$ and the consistent answer is again a
subpopulation of $D$ that the histograms of
Section~\ref{subsec:prelim-histograms} can summarize.

\paragraph{The induced pair.}
Each of the two settings supplies the framework with a pair of outputs over the
same database,
\[
(Q(D),Q'(D))
\qquad\text{and}\qquad
(Q(D),Q_\Sigma(D)) ,
\]
for refinement and for CQA respectively.
The framework requires only that each element of the pair is
a subset of $D$, and it takes the restriction of an output $T\subseteq D$ to a view
$p$ to be $T\cap p(D)$, which coincides with $Q(p(D))$ by
Proposition~\ref{prop:view-restriction} whenever $T=Q(D)$ for a selection query
$Q$. Our goal is not to compute refinements, repairs or consistent answers. The
mechanism producing the second element is treated as a black box, and the task is
to explain the difference between the two elements while preserving privacy.

\subsection{Histograms and Distributions}
\label{subsec:prelim-histograms}

Our explanations are based on distributions of attributes over query outputs and
predicate views.
Both are sets of tuples, so we define the basic statistics for an arbitrary finite
set $T$ of tuples over $\Attr(D)$; the cases of interest are $T=D$, $T=p(D)$, and
$T=Q_i(p(D))$.

\begin{definition}[Histogram]
\label{def:histogram}
Given a finite set of tuples $T$ and an attribute $A \in \Attr(D)$, the
histogram of $A$ over $T$ is the function
\[
H_{T,A}: \Dom(A)\rightarrow \mathbb{N},
\qquad
H_{T,A}(v) = |\{t\in T \mid t[A]=v\}| ,
\]
so that $H_{T,A}(v)$ is the number of tuples in $T$ whose value for $A$ equals $v$.
\end{definition}

\begin{definition}[Empirical distribution]
\label{def:empirical-distribution}
For a non-empty finite set of tuples $T$ and an attribute $A$, the empirical
distribution of $A$ over $T$ is
\[
P_{T,A}(v) = \frac{H_{T,A}(v)}{|T|},
\qquad v\in\Dom(A),
\]
equivalently $P_{T,A}(v)=\Pr\left(t[A]=v \mid t\in T\right)$ for $t$ drawn uniformly
from $T$. It is undefined for $T=\emptyset$;
\end{definition}

In the standard statistical sense, a marginal distribution is obtained by
projecting a joint distribution onto a subset of attributes. In this
paper, we use one-way marginals, namely marginals over a single
attribute.
For a query output $Q(D)$ and an attribute $A$, the one-way marginal of
$A$ is $P_{Q(D),A}$.
These one-way marginals are the basic objects used in our explanations, since
comparing $H_{Q_1(D),A}$ and $H_{Q_2(D),A}$, or their normalized counterparts,
summarizes how the two query outputs differ with respect to attribute $A$.
Such distributional summaries are also widely used in privacy-preserving data
analysis and release mechanisms~\cite{TheAlgFoundDP2014}.

For ease of presentation, we assume that the attribute values used in
histograms and predicate views are categorical. Numerical attributes are
incorporated by rendering them categorical, e.g., by bucketizing their
domains into ranges, a common practice when presenting aggregate
results. Thus, for categorical attributes the histogram bins are the
attribute values, while for numerical attributes the bins correspond to
the chosen ranges.
In either case the bins are fixed in advance from the public schema and domain
metadata and never from the data at hand, so that the set of bins carries no
information about $D$; we write $\Dom(D[A])$ for this public set of bins of
attribute $A$ and use it as the index set of every histogram over $A$.

\alon{Done: formal Bucketing definition removed and replaced with
an explanation and motivation for using buckets for numerical
attributes.}

\alon{Done: unified with the statistics discussion above.}

To quantify semantic differences between query outputs, we use statistical distance
measures between distributions.

\begin{definition}[Total Variation Distance]
\label{def:tvd}
Given two distributions $P_1$ and $P_2$ over the same domain $\Omega$, the Total
Variation Distance (TVD) between them is
\[
\TVD(P_1,P_2)
=
\frac{1}{2}
\sum_{x\in\Omega}
|P_1(x)-P_2(x)| \; \in [0,1].
\]
\end{definition}

TVD is commonly used to measure similarity between distributions and is used
throughout our framework for view filtering and representativeness analysis.
We reserve the symbols $Q$, $Q_1$ and $Q_2$ for queries and use $P_1,P_2$ for
generic distributions.

\subsection{Explanation Quality Measures}
\label{subsec:prelim-quality}

The criteria our framework optimizes, namely \emph{interestingness},
\emph{coverage}, \emph{representativeness} and \emph{diversity}, are not new to
this paper.
Each has an established meaning in data mining, data exploration and query
explanation.
Our contribution is to instantiate them for a \emph{pair} of query outputs over the
same database, and to replace them by low-sensitivity counterparts that can be
optimized under differential privacy (Section~\ref{sec:solution}).
We review the classical forms here; the specific definitions used in this paper
appear in Section~\ref{sec:explanation-quality}.

\paragraph{Interestingness.}
A pattern is interesting when it deviates from what a reference model would lead one
to expect; this is the \emph{unexpectedness} view of interestingness introduced by
Silberschatz and Tuzhilin~\cite{SilberschatzT96} and surveyed, together with the
many objective measures instantiating it, by Geng and
Hamilton~\cite{GengHamilton06} and by Tan et al.~\cite{TanKS02}.
In data management the same idea appears as deviation from an expected value in the
discovery-driven exploration of OLAP cubes~\cite{Sarawagi98}, and as
deviation between distributions in visualization recommendation, where SeeDB scores a
candidate view by the distance between the distribution it induces on a target
subset and the distribution it induces on a reference subset~\cite{SeeDB}, and
FEDEX scores an exploration step by how much it changes the distribution of an
attribute~\cite{FEDEX}.
A complementary line measures the \emph{influence} of tuples or predicates on an
aggregate answer rather than the movement of a distribution~\cite{wu2013,RoyS14}.
Our interestingness (Section~\ref{sec:explanation-quality}) is of the
distributional kind, with the two compared populations being the two query outputs
$Q_1(p(D))$ and $Q_2(p(D))$ rather than a subset and its complement.


\paragraph{Coverage.}
A deviation observed on a handful of tuples is not evidence, and every explanation
framework therefore pairs a deviation measure with a mass measure.
The oldest form is the \emph{support} of association-rule
mining~\cite{AgrawalIS93}; the same role is played by the support threshold that
MacroBase applies to its risk-ratio explanations~\cite{MacroBase}, by the support
predicate of the \textsc{Diff} operator~\cite{DIFF}, by the \emph{coverage} of an
anchor rule in model explanation~\cite{Anchors}, and by the size budget of
explanation tables~\cite{ExplanationTables}.
Coverage is also studied as a first-class property of a dataset in fairness-aware
data preparation~\cite{AsudehJJ19}.
\added{Our coverage differs from all of these in counting the evidence of \emph{both}
outputs jointly.
Since the object being explained is a comparison, a region is useful when the two
outputs together place enough tuples in it, including when only one of them does:
a region that one output fills and the other leaves empty is where a refinement or
repair adds or removes part of the population wholesale, the largest difference a
region can show.
This is what motivates the form of coverage in Section~\ref{sec:explanation-quality}.}

\paragraph{Representativeness.}
Subgroup discovery makes explicit the tension between how unusual a subgroup is and
how general it is, and resolves it with quality functions that multiply the two,
the Klösgen family $n^{a}(p-p_0)$ and weighted relative accuracy being the canonical
examples~\cite{Wrobel97,HerreraCGJ11}.
Our representativeness plays the same role with the direction of the requirement
reversed.
Because the deviation we report is the \emph{change} inside a region, the
region itself should be typical of the database rather than exceptional, so a view
is scored by the closeness of its attribute distributions to the global ones.
This guards against explanations that are dramatic only because they describe an
idiosyncratic corner of the data.

\paragraph{Diversity.}
When several results are returned together, their marginal value depends on what has
already been shown; maximal marginal relevance~\cite{MMR} is the classical
formulation and the axiomatic treatment of Gollapudi and Sharma~\cite{GollapudiS09}
and the survey of Drosou and Pitoura~\cite{DrosouP10} develop it for query and
search results.
The same principle is applied to sets of explanations in data
exploration~\cite{FEDEX,ICDE23}.
Diversity objectives of this shape are typically submodular, and are optimized
greedily~\cite{LinBilmes11}; the resulting guarantee is recalled below.

\paragraph{Set-level selection.}
Because our framework releases $k$ explanations rather than one, its objective is a
set function, and the structural property that makes greedy selection defensible is
submodularity.

\begin{definition}[Monotone submodular set function]
\label{def:submodular}
Let $C$ be a finite ground set. A function $f:2^{C}\rightarrow\mathbb{R}_{\ge 0}$ is
\emph{monotone} if $f(S)\le f(T)$ whenever $S\subseteq T\subseteq C$, and
\emph{submodular} if for all $S\subseteq T\subseteq C$ and every $c\in C\setminus T$,
\[
f(S\cup\{c\})-f(S) \;\ge\; f(T\cup\{c\})-f(T).
\]
\end{definition}

Maximizing a monotone submodular function subject to a cardinality constraint is
NP-hard in general, but the greedy rule that repeatedly adds the element of largest
marginal gain achieves a $(1-1/e)$ approximation~\cite{NemhauserWF78}, which is the
best possible in polynomial time under standard
assumptions~\cite{Feige98}.
Appendix~\ref{app:greedy} shows that the particular set objective
used here admits a stronger guarantee.

\paragraph{Quality measures under differential privacy.}
None of the measures above is designed with privacy in mind, and each is a ratio or
a product of counts whose value can be moved substantially by a single tuple.
Making them usable inside a private selection mechanism is therefore a design
problem in its own right, addressed for aggregate-query explanations by
DPXPlain~\cite{TaoGMR22} and for cluster explanations by
DPClustX~\cite{DPClustX}, both of which replace their natural quality measures by
low-sensitivity surrogates.
Section~\ref{subsec:ideal-sensitivity} quantifies the difficulty in our setting and
Section~\ref{sec:explanation-quality} constructs the measures we use.

\subsection{Differential Privacy}
\label{subsec:prelim-dp}
Because explanations are derived from sensitive data, our framework relies on
differential privacy to ensure that released explanations reveal little about any
particular individual~\cite{TheAlgFoundDP2014}. A differentially private mechanism
behaves almost the same on two databases that differ in one tuple, so whether a
single individual is included in the database has only a limited effect on the
distribution of the released output. The parameter $\epsilon$ controls the privacy
loss, and smaller values give stronger privacy at the cost of more noise.

\begin{definition}[Neighboring Datasets~\cite{DMNS2006}]
\label{def:neighboring}
Two datasets $D,D'$ are neighboring, written $D\sim D'$, if one is obtained
from the other by adding or removing a single tuple, that is, if
$D'=D\cup\{t\}$ or $D'=D\setminus\{t\}$ for some~$t$.
\end{definition}

We adopt this \emph{unbounded} form of the relation throughout, and all sensitivity
statements in Sections~\ref{sec:model} and~\ref{sec:solution} are with respect to
it.

\begin{definition}[Differential Privacy~\cite{DMNS2006,TheAlgFoundDP2014}]
A randomized mechanism $\mathcal{M}$ satisfies $\epsilon$-differential privacy if for
every pair of neighboring datasets $D\sim D'$ and every subset of outputs $S$,
\[
Pr[\mathcal{M}(D)\in S]
\leq
\exp(\epsilon)
Pr[\mathcal{M}(D')\in S].
\]
\end{definition}

The noise a mechanism must add is calibrated to the sensitivity of the quantity it
releases.

\begin{definition}[Global Sensitivity~\cite{TheAlgFoundDP2014}]
\label{def:global-sensitivity}
The global sensitivity of a function $f:\mathcal{D}\rightarrow\mathbb{R}^d$ is
\[
\Delta f
=
\max_{D\sim D'}
\|f(D)-f(D')\|_1 .
\]
\end{definition}

\begin{definition}[Laplace Mechanism~\cite{DMNS2006,TheAlgFoundDP2014}]
\label{def:laplace-mechanism}
Given a function $f:\mathcal{D}\rightarrow\mathbb{R}^d$ with global sensitivity
$\Delta f$, the Laplace mechanism releases
\[
f(D)+\left(Y_1,\ldots,Y_d\right),
\qquad
Y_1,\ldots,Y_d \stackrel{\text{i.i.d.}}{\sim} \Lap(\Delta f/\epsilon),
\]
where $\Lap(b)$ denotes the zero-mean Laplace distribution with scale $b$, whose
density is $\frac{1}{2b}\exp(-|x|/b)$.
\end{definition}

Thus, functions with lower sensitivity require less noise for the same
privacy budget, which motivates the form of the measures used in
our algorithm.
\added{A quality score written as a ratio of counts is the awkward case, because its
sensitivity is dictated by the possibility of a vanishing denominator rather than
by the candidates the mechanism actually meets.
One response is to calibrate the noise to the database at hand, as smooth
sensitivity~\cite{NissimRS07} and propose-test-release~\cite{DworkLei09} do; both
buy their tighter noise with a relaxation such as $(\epsilon,\delta)$-DP.
Section~\ref{sec:explanation-quality} takes the other route, raising every
denominator by a constant fixed in advance and independent of the data, which
keeps the guarantee in pure $\epsilon$-DP at the cost of protecting against a worst
case that may never be realised.}
When the released object is a histogram, $d$ is the number of bins and the noise is
drawn independently per bin.

\begin{definition}[Exponential Mechanism~\cite{mcsherry2007,TheAlgFoundDP2014}]
\label{def:exponential-mechanism}
Given a dataset $D$, a set of candidate outputs $Y$, a utility function
$u:Y\times D\rightarrow\mathbb{R}$, and privacy budget $\epsilon$, the exponential
mechanism selects an output $y\in Y$ with probability proportional to
\[
\exp\left(
\frac{\epsilon u(y,D)}
{2\Delta u}
\right),
\]
where $\Delta u$ is the global sensitivity of the utility function.
\end{definition}

The exponential mechanism enables private selection of high-quality explanation
candidates while preserving differential privacy.

Our algorithms select not one candidate but the top $k$, which the following mechanism does in a single pass.

\begin{definition}[One-Shot Top-$k$ Mechanism~\cite{DurfeeRogers19}]
\label{def:oneshot-topk}
Given a dataset $D$, candidates $Y$, a utility function $u$ with global sensitivity $\Delta u$, a budget $\epsilon$ and an integer $k$, the mechanism perturbs every candidate score once,
\[
\tilde{u}(y)=u(y,D)+\mathrm{Gumbel}\!\left(\tfrac{2k\,\Delta u}{\epsilon}\right),
\]
where $\mathrm{Gumbel}(b)$ is the zero-location Gumbel distribution with scale $b$
and the draws are independent across candidates,
and returns the $k$ candidates with the largest $\tilde{u}$, in that order.
\end{definition}

\begin{theorem}[\cite{DurfeeRogers19}]
\label{thm:oneshot-topk}
The one-shot top-$k$ mechanism satisfies $\epsilon$-differential privacy.
\end{theorem}

The guarantee follows because the mechanism is distributionally identical to applying the exponential mechanism $k$ times without replacement, each with budget $\epsilon/k$, so sequential composition (Theorem~\ref{thm:sequential-composition}) applies; the practical gain is that the scores are perturbed once rather than recomputed in each of the $k$ rounds.
The same mechanism is used for private top-$k$ selection by DPXPlain~\cite{TaoGMR22} and by the clustering-explanation framework of~\cite{DPClustX}.


Differential privacy satisfies several important composition properties, enabling
complex privacy-preserving analytical workflows constructed from multiple private
mechanisms \cite{mcsherry2009,DBLP:series/synthesis/2016Li,bun2016}.
These properties are central to the privacy analysis of our algorithms, which combine
private view selection, private candidate selection, and noisy histogram release.

\begin{theorem}[Sequential Composition~\cite{TheAlgFoundDP2014,DBLP:series/synthesis/2016Li}]
\label{thm:sequential-composition}
Let $\mathcal{M}_1,\ldots,\mathcal{M}_k$ be randomized mechanisms such that
$\mathcal{M}_i$ satisfies $\epsilon_i$-differential privacy for every
$i \in \{1,\ldots,k\}$.
Then the mechanism that releases the joint output
\[
\mathcal{M}(D)
=
\left(
\mathcal{M}_1(D),
\ldots,
\mathcal{M}_k(D)
\right)
\]
satisfies
\[
\left(\sum_{i=1}^{k}\epsilon_i\right)
\text{-differential privacy}.
\]
The statement holds in the \emph{adaptive} setting as well.
If the choice of $\mathcal{M}_i$ may depend on the outputs of $\mathcal{M}_1,\dots,\mathcal{M}_{i-1}$, the same bound $\sum_i\epsilon_i$ applies.
\end{theorem}

We rely on the adaptive form, since the utility optimised in each round of Algorithm~\ref{alg:topk} depends on the explanations released in earlier rounds.

Sequential composition allows us to account for the total privacy loss of an algorithm
that invokes multiple differentially private subroutines by summing their individual
privacy budgets.

\begin{theorem}[Post-Processing~\cite{TheAlgFoundDP2014,DBLP:series/synthesis/2016Li}]
\label{thm:post-processing}
Let $\mathcal{M}$ be an $\epsilon$-differentially private mechanism,
and let $f$ be any randomized mapping that does not access the private database
except through the output of $\mathcal{M}$.
Then the mechanism
\[
f(\mathcal{M}(D))
\]
also satisfies $\epsilon$-differential privacy.
\end{theorem}

The post-processing property implies that any deterministic or randomized computation
performed only on already-private outputs does not incur additional privacy loss.
In particular, the constants $\tau$ and $\gamma$ of
Equations~\eqref{eq:tau-rule} and~\eqref{eq:gamma-rule}, being functions of public quantities alone,
consume no privacy budget.
